\section{Listings 代码块}

\begin{lstlisting}[language=Python, style=mystyle]
def kmeans_clustering(data, k, max_iter=100):
    n_samples, n_features = data.shape
    centroids = data[np.random.choice(n_samples, k, replace=False)]
    
    for _ in range(max_iter):
        distances = np.sqrt(((data - centroids[:, np.newaxis])**2).sum(axis=2))
        labels = np.argmin(distances, axis=0)
        
        new_centroids = np.array([data[labels == i].mean(axis=0) for i in range(k)])
        
        if np.allclose(centroids, new_centroids):
            break
        centroids = new_centroids
    
    return labels, centroids
\end{lstlisting}

\begin{lstlisting}[language=HTML, style=mystyle]
<!DOCTYPE html>
<html lang="zh-CN">
<head>
    <meta charset="UTF-8">
    <title>示例页面</title>
    <link rel="stylesheet" href="styles.css">
</head>
<body>
    <div class="container">
        <h1>欢迎使用 Yan Book 模板</h1>
        <p>这是一个示例段落。</p>
        <ul>
            <li>列表项 1</li>
            <li>列表项 2</li>
        </ul>
    </div>
</body>
</html>
\end{lstlisting}

\begin{lstlisting}[style=mystyle]
.container {
    max-width: 1200px;
    margin: 0 auto;
    padding: 20px;
}

.title {
    font-size: 24px;
    color: #333;
    margin-bottom: 16px;
}

.button {
    display: inline-block;
    padding: 8px 16px;
    background-color: #007bff;
    color: white;
    border: none;
    border-radius: 4px;
    cursor: pointer;
    transition: background-color 0.3s;
}

.button:hover {
    background-color: #0056b3;
}
\end{lstlisting}